\documentclass{proc}
\usepackage{blindtext}
\usepackage{titlesec}
\usepackage{graphicx}
%\usepackage{underscore}
\usepackage{float}
\usepackage{tabularx}
\usepackage{subfigure}
\usepackage{hyperref}
\usepackage{pdfpages}
\usepackage{setspace}

\bibliographystyle{plain} 

\title{%
	Crowd Counting on Edge Devices \\
	\large DA642E Project}
\author{Adrian Jon Berlovan\and Athanasios Smilkos \and Vrysiida Charizi}
\date{\today}

\begin{document}

\maketitle

\clearpage

\begin{abstract}
In the current day and age, our right to privacy is being infringed more and more often. The goal of this project is to help people regain that right, especially in public settings. Often time cameras located in public areas need to send data back to the cloud in order to do inference or segmentation, not caring about privacy. In order to do image segmentation in a privacy-friendly way, we will be investigating running such machine learning (ML) algorithms directly on edge devices. This subject has been often tackled in literature, however usually using some form of hybrid edge, server and federated learning. 

We have combined several previously existing datasets into a new one, with a smaller selection of persons per image, which is more representative of small gatherings in public places where people might be easily identifiable. The YOLO26 pre-trained network has been re-trained on this dataset and employed quantization in order to be able to run the model faster on lower end hardware.
\end{abstract}

\section{Introduction} % Describe the problem you are working on, data source(s), and motivation
\section{Related Work} % Discuss published works that relate to your project

The machine learning model we are basing our project on is YOLO\cite{YOLO}, a pre-trained model, currently on version 26. It comes in various sizes, from n (the smallest, at 5 million parameters) all the way to x (the largerst, at 100 million parameters). Much other work has been done in this domain\cite{UAV}, however the models are usually not relased, and as such the claims are tougher to verify. Many of the related research does crowd counting without segmentation, which is good for very large gatherings but might prove less applicable to smaller groups of people.

\section{Method} % Describe the data you are working with for your project. The pre-processing steps, and AI method(s) employed to solve the problem/gain insights.
\section{Results} %  Describe the experiments that you performed to demonstrate that your approach solves the problem. Discuss the results and compare them with related works.
We have tried on ESP32, however that did not work at all for reasons that we will describe below:
\section{Conclusions} % (and ideas for future work)

\bibliography{bibliography}

\end{document}

% In the current day and age, we see the right to privacy being disregarded more and more often. In our project, we aim to help people regain that right. We will be investigating running crowd detection machine learning (ML) algorithms directly on edge devices. To the best of our knowledge, this subject has not been tackled that often in the literature, or has used some form of hybrid edge, server, and federated learning \cite{UAV}.

%\section{Application area}

% Our main concern relates to counting crowds in a privacy preserving way. As a project example, we can use Skånetrafiken ticket control. Instead of having controllers blindly board buses and bother people, we propose the use of an edge camera that can gather crowd statistics offline. By comparing that with the amount of scanned tickets, we can infer which lines are more prone to fraud. In this way, problematic lines can be isolated and investigated further by controllers.

% To that end, we will compare the performance of different edge devices (Raspberry Pis, Laptops, and other similarly weak hardware) in running different state of the art architectures \cite{RepSFNet} \cite{santos2025wireless} \cite{oneshot}. In order to run the models, we will attempt to minimize the neural networks, by employing various techniques such as pruning, quantization or knowledge distillation.

% \section{Data}

% For our benchmarks, we will be evaluating and comparing different datasets, and finally choose one of them \cite{CrowdHuman} \cite{COCO} \cite{KITTI}. Since we are interested in a smaller crowd of people (between 5-20), we will have to analyze the dataset and extract only the samples that fit our criteria. With that, we will explore retraining the existing models on our new data.

% \section{Results}

% To evaluate whether the algorithm is applicable on the tested platform, we will be conducting a series of experiments:
% \begin{itemize}
% \item Timing - will it run in a reasonable amount of time?
% \item Scale - can it fit on the device?
% \item Accuracy - does it count the number of people correctly?
% \end{itemize}

% After running the tests and collecting the necessary results, we will compile it all in graph form, for easier visualization. 

